\chapter{Kết quả đạt được, hạn chế và đề xuất}

Quá trình thực tập tại phòng thí nghiệm kiểm định chất lượng sản phẩm của Công ty TNHH Quốc Tế Khánh Sinh vận dụng trực tiếp các kiến thức Hóa phân tích vào môi trường sản xuất thực tế. Việc tham gia lấy mẫu, chuẩn bị mẫu, pha hóa chất, thực hiện các phép phân tích và xử lý số liệu là cơ hội tiếp cận các chỉ tiêu kiểm nghiệm quan trọng của phân bón như nitơ tổng số, phospho hữu hiệu, kali hữu hiệu, carbon hữu cơ, chất hữu cơ và độ ẩm. Đồng thời, quá trình thực tập còn giúp làm quen với việc sử dụng các thiết bị như cân phân tích, hệ thống chưng cất, máy quang kế ngọn lửa, tủ sấy và các dụng cụ phân tích thể tích.

Bên cạnh kiến thức chuyên môn, nhiều kỹ năng nghề nghiệp cũng được hình thành và củng cố, bao gồm thao tác phòng thí nghiệm chính xác, quản lý mẫu, ghi nhận số liệu, tính toán kết quả, đánh giá độ lặp và đối chiếu kết quả thành phẩm với chỉ tiêu đăng ký. Qua đó, mối liên hệ giữa lý thuyết Hóa phân tích với thực tế kiểm soát chất lượng, đặc biệt là vai trò của việc chuẩn hóa thao tác, kiểm soát sai số và đảm bảo khả năng truy xuất của kết quả được thể hiện rõ hơn.

Trong quá trình thực tập vẫn ghi nhận một số hạn chế như một số phép thử còn phụ thuộc nhiều vào thao tác thủ công, một số mẫu có độ lặp chưa đạt yêu cầu hoặc kết quả thành phẩm khác với mức công bố. Do thời gian thực tập còn hạn chế, các trường hợp này chưa được phân tích lại đầy đủ để xác định nguyên nhân chính xác. Vì vậy, một số giải pháp được đề xuất gồm chuẩn hóa thao tác, tăng cường kiểm tra thiết bị và đường chuẩn, sử dụng mẫu QC định kỳ, thực hiện lại các mẫu không đạt và từng bước số hóa việc ghi nhận, tính toán và lưu trữ dữ liệu. Nhìn chung, kỳ thực tập đã giúp củng cố nền tảng chuyên môn, nâng cao kỹ năng thực hành và hình thành tư duy kiểm soát chất lượng có hệ thống trong lĩnh vực phân tích phân bón.

\section{Những kiến thức và kinh nghiệm đạt được trong quá trình thực tập}

Quá trình thực tập tại phòng thí nghiệm kiểm định chất lượng sản phẩm của Công ty TNHH Quốc Tế Khánh Sinh đã tạo điều kiện để vận dụng các kiến thức lý thuyết vào hoạt động kiểm nghiệm thực tế tại doanh nghiệp sản xuất phân bón. 
Việc trực tiếp quan sát và tham gia vào các công việc của phòng thí nghiệm là cơ hội tiếp cận toàn bộ chuỗi kiểm soát chất lượng từ lấy mẫu nguyên liệu, bán thành phẩm và thành phẩm đến chuẩn bị mẫu, tiến hành phân tích, xử lý số liệu và đánh giá kết quả.

Một trong những kiến thức quan trọng thu nhận được là cách tổ chức quy trình lấy và quản lý mẫu. 
Mẫu kiểm nghiệm không chỉ cần được lấy đúng vị trí và có tính đại diện mà còn phải được mã hóa, ghi nhận đầy đủ thông tin và lưu trữ phù hợp để đảm bảo khả năng truy xuất. 
Đối với mẫu thành phẩm, kết quả kiểm nghiệm còn được đối chiếu với chỉ tiêu chất lượng đăng ký để đánh giá mức độ phù hợp trước khi sản phẩm được đưa ra thị trường.

Bên cạnh đó, quá trình thực tập giúp củng cố kiến thức về các phương pháp phân tích được sử dụng trong kiểm nghiệm phân bón. 
Một số chỉ tiêu được tiếp cận và thực hành gồm nitơ tổng số, phospho hữu hiệu, kali hữu hiệu, carbon hữu cơ/chất hữu cơ và độ ẩm. 
Mỗi chỉ tiêu sử dụng một nguyên lý phân tích khác nhau, điều này là do việc lựa chọn phương pháp phụ thuộc vào bản chất hóa học của chất cần xác định.

Kết quả phân tích phụ thuộc vào nhiều yếu tố nhỏ trong quá trình thao tác như độ đồng nhất của mẫu, khối lượng mẫu, độ chính xác khi hút pipet, định mức, thời gian phản ứng, xác định điểm cuối chuẩn độ, điều kiện vận hành thiết bị và cách ghi nhận số liệu.

Bên cạnh kiến thức chuyên môn, quá trình thực tập giúp hình thành nhiều kỹ năng cần thiết đối với công việc tại phòng kiểm nghiệm. 
Quá trình thực tập còn rèn luyện tính cẩn thận, kiên nhẫn và tinh thần trách nhiệm, đặc biệt trong các thao tác cân, pha dung dịch, chuẩn độ và ghi nhận số liệu, bởi những sai lệch nhỏ trong quá trình thực hiện có thể ảnh hưởng trực tiếp đến kết quả kiểm nghiệm. 
Kỹ năng làm việc nhóm và phối hợp công việc cũng được cải thiện thông qua việc trao đổi với nhân viên phòng thí nghiệm, Phòng Kỹ thuật và Bộ phận Sản xuất trong quá trình lấy mẫu, kiểm tra và xử lý kết quả. 
Ngoài ra, khi xuất hiện các kết quả bất thường hoặc độ lặp chưa đạt yêu cầu, khả năng nhận diện vấn đề, phân tích nguyên nhân và đề xuất hướng xử lý được rèn luyện thông qua việc xem xét lại các yếu tố như tính đồng nhất của mẫu, thao tác phân tích, hóa chất, thiết bị và điều kiện thử nghiệm.

\section{Mối liên hệ giữa nội dung thực tập và kiến thức lý thuyết}

Nội dung thực tập có mối liên hệ trực tiếp với các kiến thức được học trong môn Hóa phân tích (Analytical Chemistry). Nếu trong quá trình học, các phương pháp thường được tiếp cận thông qua nguyên lý, phương trình phản ứng và các bài toán tính toán thì tại phòng thí nghiệm, những kiến thức này được đặt trong một quy trình phân tích hoàn chỉnh với mẫu thực tế.

Các nội dung về phân tích thể tích được ứng dụng rõ rệt trong phép xác định nitơ tổng số và carbon hữu cơ. Những kiến thức về phản ứng acid--base, phản ứng oxy hóa--khử, chuẩn độ ngược, dung dịch chuẩn, chỉ thị và điểm kết thúc được sử dụng trực tiếp trong quá trình phân tích. Qua thực tế có thể nhận thấy rằng việc xác định đúng điểm kết thúc chuẩn độ và chuẩn bị chính xác nồng độ dung dịch chuẩn ảnh hưởng đáng kể đến độ tin cậy của kết quả.

Kiến thức về phân tích công cụ được vận dụng trong phép xác định kali bằng quang kế ngọn lửa. Quá trình xây dựng đường chuẩn giúp liên hệ trực tiếp giữa lý thuyết về hiệu chuẩn thiết bị và thực tế định lượng mẫu chưa biết. Việc xây dựng phương trình hồi quy, đánh giá hệ số $R^2$, xác định nồng độ từ tín hiệu thiết bị và sử dụng hệ số pha loãng cho thấy vai trò của xử lý số liệu trong phân tích định lượng.

Các kiến thức về chuẩn bị mẫu và xử lý mẫu cũng được thể hiện xuyên suốt quá trình thực tập. Một kết quả phân tích chính xác không chỉ phụ thuộc vào thiết bị hay phép chuẩn độ mà bắt đầu từ việc lấy được mẫu đại diện và xử lý mẫu phù hợp. Điều này giúp làm rõ một nội dung quan trọng của Hóa phân tích: sai số có thể phát sinh ở mọi giai đoạn của quy trình và kết quả cuối cùng chỉ có ý nghĩa khi toàn bộ chuỗi phân tích được kiểm soát.

Ngoài ra, việc tính RPD giữa các lần thử giúp liên hệ với các khái niệm độ chụm, độ lặp và kiểm soát chất lượng kết quả phân tích. Quá trình thực tập vì vậy đã bổ sung phần ứng dụng thực tế cho những kiến thức Hóa phân tích đã được học, đồng thời giúp hiểu rõ hơn ý nghĩa của từng bước trong một phương pháp thử tiêu chuẩn.

\section{Một số hạn chế ghi nhận trong quá trình thực tập}

Trong phạm vi thời gian thực tập, một số hạn chế được ghi nhận chủ yếu liên quan đến quá trình kiểm nghiệm và quản lý kết quả. Những nhận xét này được đưa ra dựa trên phạm vi công việc đã tiếp cận và không đại diện cho toàn bộ hoạt động quản lý chất lượng của công ty.

Thứ nhất, một số phép phân tích vẫn phụ thuộc tương đối nhiều vào thao tác thủ công của người thực hiện, đặc biệt ở các bước cân, pha dung dịch, chuẩn độ, nhận biết điểm cuối và nhập số liệu. Điều này có thể làm tăng sai số giữa các lần thử nếu thao tác chưa hoàn toàn đồng nhất.

Thứ hai, trong số liệu thực tế vẫn xuất hiện một số trường hợp RPD giữa hai lần thử vượt giới hạn chấp nhận, đặc biệt đối với phép xác định kali. Do thời gian thực tập có hạn, các mẫu này chưa được thực hiện lại đầy đủ để xác định chính xác nguyên nhân.

Thứ ba, kết quả thành phẩm cho thấy một số mẫu có chỉ tiêu K$_2$O, P$_2$O$_5$ hoặc độ ẩm chưa phù hợp với mức đăng ký. Các sai lệch này có thể xuất phát từ nhiều nguồn khác nhau, từ quá trình phân tích đến chất lượng nguyên liệu, tỷ lệ phối trộn, độ đồng nhất của hỗn hợp hoặc điều kiện sản xuất và bảo quản. Trong thời gian thực tập chưa đủ điều kiện để thực hiện một nghiên cứu nguyên nhân đầy đủ cho từng trường hợp.

Thứ tư, việc tổng hợp và xử lý số liệu trên các bảng biểu có thể làm tăng khối lượng thao tác thủ công, đặc biệt khi số lượng mẫu và chỉ tiêu kiểm nghiệm tăng. Việc nhập liệu nhiều lần cũng tạo nguy cơ phát sinh sai sót trong quá trình chuyển số liệu từ phiếu thử nghiệm sang bảng tổng hợp.

Cuối cùng, số lần thử song song trong kiểm nghiệm thường ở mức hai lần, phù hợp với công tác kiểm tra thường quy nhưng còn hạn chế nếu muốn đánh giá sâu các đặc tính của phương pháp như độ chụm trung gian, độ đúng, độ tái lập hoặc độ ổn định dài hạn.

\section{Đề xuất giải pháp}

Để nâng cao độ tin cậy và hiệu quả của hoạt động kiểm nghiệm, trước hết có thể tăng cường chuẩn hóa thao tác thông qua các hướng dẫn thao tác chi tiết cho từng phép thử. Các bước có ảnh hưởng lớn đến kết quả như lượng mẫu, thời gian phản ứng, tốc độ chuẩn độ, cách xác định điểm kết thúc và điều kiện vận hành thiết bị nên được quy định rõ nhằm giảm sự khác biệt giữa những người thực hiện.

Đối với các mẫu có độ lặp không đạt, nên thiết lập quy trình xử lý cụ thể gồm kiểm tra lại số liệu tính toán, kiểm tra dụng cụ và thiết bị, đồng nhất lại mẫu và thực hiện lại phép thử từ bước cân mẫu. Nếu kết quả tiếp tục không đạt, cần ghi nhận như một trường hợp không phù hợp và tiến hành điều tra nguyên nhân.

Đối với các mẫu thành phẩm có kết quả khác với mức đăng ký, cần phân biệt nguyên nhân do quá trình phân tích với nguyên nhân do sản xuất. Sau khi xác nhận kết quả kiểm nghiệm bằng phép thử lặp lại, có thể truy ngược về kết quả nguyên liệu, phiếu phối liệu, lượng nguyên liệu được cấp vào dây chuyền, thời gian phối trộn và điều kiện bảo quản để xác định vị trí phát sinh sai lệch.

Công ty cũng có thể tăng cường sử dụng mẫu kiểm soát chất lượng (QC) được phân tích định kỳ. Việc theo dõi cùng một mẫu QC trong thời gian dài và xây dựng biểu đồ kiểm soát sẽ giúp nhận biết sớm xu hướng dịch chuyển của phương pháp hoặc thiết bị trước khi sai lệch ảnh hưởng đến kết quả mẫu thực tế.

Đối với máy quang kế ngọn lửa và các thiết bị phân tích khác, nên duy trì việc kiểm tra đường chuẩn, mẫu trắng và dung dịch kiểm soát theo tần suất quy định. Kết quả kiểm tra thiết bị cần được lưu cùng hồ sơ phân tích để hỗ trợ việc truy xuất khi xuất hiện kết quả bất thường.

Về quản lý số liệu, có thể từng bước số hóa quá trình ghi nhận và tính toán kết quả. Các công thức tính OC, OM, N, P$_2$O$_5$, K$_2$O, độ ẩm, giá trị trung bình, RPD và đối chiếu giới hạn chất lượng có thể được xây dựng thành các bảng tính có kiểm soát hoặc hệ thống quản lý dữ liệu phòng thí nghiệm. Việc tự động hóa các bước tính toán giúp giảm lỗi nhập liệu và tạo thuận lợi cho việc truy xuất lịch sử kết quả theo mã mẫu hoặc lô sản xuất.

Cuối cùng, nếu có điều kiện, phòng thí nghiệm có thể thực hiện các chương trình đánh giá phương pháp định kỳ với số lượng phép lặp lớn hơn, mẫu chuẩn hoặc mẫu thêm chuẩn để đánh giá đồng thời độ lặp, độ chụm trung gian và độ đúng. Điều này sẽ giúp dữ liệu kiểm nghiệm không chỉ phục vụ quyết định đạt hay không đạt của từng lô sản phẩm mà còn cung cấp thông tin về hiệu năng lâu dài của chính phương pháp phân tích.
